\section{Comparison with fundamental-constant adjustments}
\label{sec:comparacion}

External comparison evaluates the consequences of expressions whose initial states, regions, coefficients, and carries have been determined by the preceding construction. For each reference \(x_{\rm ref}\), with standard uncertainty \(u(x_{\rm ref})\), the normalized residual is reported:
\[
 z=\frac{x_{\rm calc}-x_{\rm ref}}{u(x_{\rm ref})}.
\]
This ratio expresses the difference from each tabulated centre in units of its standard uncertainty. A probabilistic interpretation requires a specified statistical model; joint comparison of the CODATA adjustments additionally requires their correlations and the independence conditions being used.

The expressions used in this comparison give
\begin{align}
 \alpha&=0.007297352569283800997285105472380662999\ldots,\notag\\
 \frac{m_ec^2}{\mathrm{MeV}}&=0.510998950690000808608\ldots.
 \label{eq:valores-comparacion}
\end{align}
The second value corresponds to the corrected electronic expression, with the complete \(\Delta_4\) factor in the basal normalization and in the return exponent. Composition of the basal term with return memory produces this value before the tabulated mass is introduced as a comparison reference.

\subsection{Exactness, uncertainty, and dimensional realization}

The comparison connects three quantities with distinct functions. Algebraic exactness concerns identities proved within their domains. Numerical evaluation error is controlled through computational precision. Experimental uncertainty quantifies knowledge of the measured quantity under the adjustment used and retains that meaning when the decimal evaluation is extended. The additional digits in \eqref{eq:valores-comparacion} describe the mathematical expression; experimental information is represented by the centres, uncertainties, and correlations of the adjustment.

The dimensional realization in Section~\ref{sec:electron} specifies the map from the electronic character to energy and mass, with its scale and normalization. This map determines the quantity and units on which the comparison acts. The structural construction establishes the evaluated value and its relations; metrological comparison examines its correspondence with the physical quantity thus identified. The causal order retains APP, TRIT, and TPK as antecedents of the evaluated expressions, followed by dimensional realization and the reference values.

\subsection{CODATA 2018}

The 2018 adjustment contains the following central values and uncertainties\cite{codata2018}:
\begin{center}
\begin{tabular}{@{}lcccr@{}}\toprule
Quantity&HMT value&Reference&\shortstack{Standard\\uncertainty}&\(z\)\\\midrule
\(\alpha\)&\(0.0072973525692838\ldots\)&\(0.0072973525693\)&\(1.1\times10^{-12}\)&\(-0.01473\)\\
\(m_ec^2/\mathrm{MeV}\)&\(0.5109989506900008\ldots\)&\(0.51099895000\)&\(1.5\times10^{-10}\)&\(+4.60001\)\\\bottomrule
\end{tabular}
\end{center}
The expression for \(\alpha\) agrees with the 2018 central value when rounded to the digits shown. The difference, approximately \(-1.62\times10^{-14}\), is about \(0.0147\) times the printed standard uncertainty. The corrected electronic expression exceeds the 2018 centre by approximately \(6.90\times10^{-10}\,\mathrm{MeV}\), a difference that persists at the stated rounding. Each row thus retains its decimal agreement or its quantified discrepancy.

\subsubsection{Comparison of the dodecaphase window with the reciprocal centre}
\label{sec:contraste-ventana-reciproco}

The finite window \(\alpha_{12}^{(0)}\) in
\eqref{eq:alpha-ventana-cero}, obtained with \(c_{13}=0\), also admits
the historical comparison with the reciprocal of the printed central
value of the inverse constant. Define
\[
 R_{2018}:=137.035999084,\qquad
 \alpha_{\mathrm{rec},2018}:=R_{2018}^{-1}.
\]
This operation evaluates the reciprocal of the decimal rational
\(R_{2018}\). This reference is distinct from the tabulated central
value \(0.0072973525693\) used in the preceding table.

\begin{center}
\begin{tabular}{@{}ll@{}}
\toprule
Expression & Arithmetic evaluation\\
\midrule
\(\alpha_{12}^{(0)}\) &
\(0.007297352569283800997285105472380663\)\\
\(\alpha_{\mathrm{rec},2018}\) &
\(0.007297352569283800997285105472380663567972560701\ldots\)\\
\(\alpha_{12}^{(0)}-\alpha_{\mathrm{rec},2018}\) &
\(\simeq-5.679725607012965\times10^{-37}\)\\
\((\alpha_{12}^{(0)})^{-1}-R_{2018}\) &
\(\simeq+1.066588006664435\times10^{-32}\)\\
\(\bigl((\alpha_{12}^{(0)})^{-1}-R_{2018}\bigr)/R_{2018}\) &
\(\simeq+7.783268730800000\times10^{-35}\)\\
\bottomrule
\end{tabular}
\end{center}

The precise decimal relation is
\[
 \alpha_{12}^{(0)}
 =10^{-36}\left\lfloor10^{36}\alpha_{\mathrm{rec},2018}\right\rfloor.
\]
The historical agreement therefore concerns truncation of the
reciprocal to thirty-six decimal places; rounding at that position
gives a final block \(664\), whereas the window ends in \(663\).
The differences in the table quantify the comparison between these
two numbers. The infinite section \(\alpha_A\) retains its own
compatible normalization and boundary, distinguished from this
finite window.

The adjustment reports the inverse constant as
\(137.035999084(21)\), with standard uncertainty
\(u(R_{2018})=2.1\times10^{-8}\)\cite{codata2018}.
The decimal digits obtained by inverting its printed centre belong
to the arithmetic evaluation of that rational; experimental uncertainty
retains the value specified by the adjustment. The preceding comparison
with the tabulated centre of \(\alpha\) retains its own reference and
normalized residual.


\subsection{CODATA 2022}

The 2022 adjustment changes the central values\cite{codata2022}:
\begin{center}
\begin{tabular}{@{}lcccr@{}}\toprule
Quantity&HMT value&Reference&\shortstack{Standard\\uncertainty}&\(z\)\\\midrule
\(\alpha\)&\(0.0072973525692838\ldots\)&\(0.0072973525643\)&\(1.1\times10^{-12}\)&\(+4.53073\)\\
\(m_ec^2/\mathrm{MeV}\)&\(0.5109989506900008\ldots\)&\(0.51099895069\)&\(1.6\times10^{-10}\)&\(+0.00000505\)\\\bottomrule
\end{tabular}
\end{center}
In this edition, the electronic expression agrees with the tabulated central value after rounding, whereas \(\alpha\) differs by approximately \(4.53\) printed standard uncertainties. Comparison of the physical accuracy of the adjustments rests on their experimental data, correlations, and assumptions. The date and decimal agreement with an expression identify properties of the comparison; assessing proximity to the physical value requires that complete experimental analysis.

\Needspace{29\baselineskip}
\subsection{Comparison of the reduced Planck constant}
\label{sec:revision-planck-contraste}

The International System in force since 20 May 2019 fixes
\(h=6.62607015\times10^{-34}\,\mathrm{J\,s}\) exactly
\cite{bipm2018}. Its reduced constant is therefore
\[
 \hbar_{\mathrm{SI}}=\frac{6.62607015}{2\pi}\,10^{-34}\,\mathrm{J\,s}
 =1.0545718176461563912\ldots\times10^{-34}\,\mathrm{J\,s}.
\]
The ellipsis denotes the expansion of an exact expression,
not an experimental uncertainty in \(h\) in the current SI.
This reference value is used here after evaluation of the
HMT readers; its function is comparative.

The article's two return coordinates are
\begin{align*}
 H_5&=1.0545718176461544696\ldots,\\
 \eta_{\mathrm{ret}}&=1.0545718169150390611\ldots.
\end{align*}
Under the identification of \(\mathcal U_S\) with \(\mathrm{J\,s}\), the
relative difference between the preceding section and
\(\hbar_{\mathrm{SI}}\) is approximately \(-1.82\times10^{-15}\).
For the subsequent section, it is approximately \(-6.93\times10^{-10}\).
The second difference contains the return effect used by the
electronic corrector; it does not equal the evaluation error of the
first section.

The precision of these comparisons allows the quantitative proximity
of the construction of \(H_5\) to be recognized while distinguishing
proximity from exact identity. In the present manuscript, the displayed
expressions do not provide an exact equality with \(h/(2\pi)\).
The result proved for \(\Sigma_H\) is its decade \(-34\);
the analytic result for \(H_5\) is the evaluation of the stated
character. The physical identification is examined in this chart and with the
explicit differences just given. The SI value does not participate
in the definition of the determinantal norm or in the regional recurrence.


